\documentclass[11pt,a4paper]{article}
\usepackage[margin=25mm]{geometry}
\usepackage{amsmath,amssymb,booktabs,longtable}
\usepackage[T1]{fontenc}
\usepackage{lmodern}
\usepackage[hidelinks]{hyperref}
\newcommand{\ind}{\mathbf{1}}
\newcommand{\avg}[1]{\sum_{k=0}^{N}\pi_k #1}
\title{Mathematical Analysis of UL Trigger-Based DPS\\
\large Detailed Derivation of the Implemented Saturated Approximation}
\author{}
\date{}
\begin{document}
\maketitle
\begin{abstract}
A station first wins contention to report an EBSR; the AP subsequently wins
separate contention to serve that report using an HC uplink exchange. We
construct a reduced state-dependent contention model and derive its
stationary probabilities, channel holding times, throughput, collision ratio
per attempt, time occupancy, and all-STA energy rewards. The equations match
\texttt{src/dps\_sim/analytical.py}. The reduced chain is an approximation,
not an exact Markov description of the packet simulator.
\end{abstract}
\tableofcontents

\section{Operation, assumptions, and units}
There are $N\geq1$ saturated STAs and one AP sharing a serialized channel.
A successful packet requires two separate contention wins:
\begin{enumerate}
\item The STA wins a TXOP, sends EBSR on LC, and leaves contention while
waiting for its report to be served.
\item The AP wins a later TXOP and serves the oldest report: ICF on LC,
switch up, ICR--DATA--ACK on HC, then switch down to LC.
\end{enumerate}
The AP is one contender whenever reports are pending. Each report requests
one packet; each grant serves one packet, FIFO. After delivery the saturated
STA immediately has a replacement packet. Backoff counters freeze during
busy events. Retries are unlimited, but the stage is capped at $m$.

We assume identical contention parameters, no PHY errors, hidden nodes,
capture, aggregation, or concurrent exchanges. EBSR and ICF have equal
lengths and rates. Switching occupies channel time and does not overlap
interframe gaps. Only the served STA switches to HC; AP energy is excluded.
These assumptions define a research abstraction, not a full IEEE implementation.

Durations use microseconds, rates use bits per microsecond (numerically Mbps),
lengths use bits, and powers use watts. Let $\gamma=10^{-6}$ seconds per
microsecond. A frame of $b$ bits at rate $R$ occupies $b/R$ microseconds
and consumes $\gamma(b/R)P$ joules at power $P$.
\begin{longtable}{@{}lll@{}}
\toprule Symbol & Meaning & Default\\\midrule
$N$ & Number of STAs & varied\\
$W,m$ & Initial window size, maximum stage & $32,5$\\
$\sigma$ & Idle slot & $50\,\mu$s\\
$s,d,\delta$ & SIFS, DIFS, propagation & $28,128,1\,\mu$s\\
$L,h$ & Payload, DATA header & $8184,400$ bits\\
$e,f$ & EBSR, ICF including PHY & $288,288$ bits\\
$r,a$ & ICR, ACK including PHY & $240,240$ bits\\
$R_C,R_H$ & Control, HC DATA rate & $1,10$ Mbps\\
$u,v$ & Switch up, switch down & $100,100\,\mu$s\\
$P_{L,i},P_{L,r},P_{L,t}$ & LC idle, RX, TX & $0.10,0.20,0.50$ W\\
$P_{H,i},P_{H,r},P_{H,t}$ & HC idle, RX, TX & $0.30,0.50,1.00$ W\\
$P_{sw}$ & Served-STA switching power & $0.40$ W\\
\bottomrule
\end{longtable}

\section{Local contention model}
\subsection{Active population}
Let $K=k\in\{0,\ldots,N\}$ count reported STAs awaiting service. Define
\begin{equation}
n_k=N-k,\qquad b_k=\ind_{\{k>0\}},\qquad M_k=n_k+b_k.
\end{equation}
There are $n_k$ eligible STAs and at most one active AP contender. In
particular, $M_0=N$ and $M_N=1$. Treating DPS as $N+1$ always-active nodes
would ignore the withdrawal of reported STAs.

\subsection{Backoff attempt probability}
For a tagged contender, approximate the collision probability by a constant
$p<1$ during a local equilibrium calculation. At stage $i$, its counter is
uniform on $\{0,\ldots,W2^i-1\}$. Consider a renewal cycle ending at its next
success. Expected stage-visit weights are
\begin{equation}
v_i=p^i\quad(0\leq i<m),\qquad v_m=\frac{p^m}{1-p}.
\end{equation}
The capped stage can be revisited indefinitely. The mean attempts per cycle
are $\sum_i v_i=1/(1-p)$. Each visit contributes one attempt and an expected
$(W2^i-1)/2$ idle countdown opportunities. Therefore the attempt probability
in the decoupled virtual-opportunity description is
\begin{align}
\tau(p)&=\frac{\sum_i v_i}{\sum_i v_i[1+(W2^i-1)/2]}\\
&=\frac{2}{1+W(1-p)\sum_i v_i2^i}.
\end{align}
Using the identity
\begin{align}
(1-p)\sum_{i=0}^{m}v_i2^i
&=(1-p)\sum_{i=0}^{m-1}(2p)^i+(2p)^m\\
&=1+p\sum_{j=0}^{m-1}(2p)^j,
\end{align}
we obtain
\begin{equation}\label{eq:tau}
\boxed{\tau(p)=\frac{2}{W+1+Wp\sum_{j=0}^{m-1}(2p)^j}}.
\end{equation}
The polynomial form is valid at $p=1/2$; for $m=0$ the sum is empty.
For state $k$, assume independent attempts by all $M_k$ contenders:
\begin{equation}\label{eq:fixed}
p_k=1-(1-\tau_k)^{M_k-1},\qquad \tau_k=\tau(p_k).
\end{equation}
A tagged attempt collides when at least one other contender transmits. When
$M_k=1$, $p_k=0$ and $\tau_k=2/(W+1)$.

This is a local equilibrium closure. The real simulator retains individual
counters, stages, and FIFO identities. A change in $k$ does not actually
reset the remaining contenders to independent equilibrium samples.

\section{Virtual events and the embedded chain}
A virtual event is one idle countdown slot or a complete busy event. Its
four mutually exclusive outcomes have probabilities
\begin{align}
P_0(k)&=(1-\tau_k)^{M_k},\\
P_E(k)&=n_k\tau_k(1-\tau_k)^{M_k-1},\\
P_A(k)&=b_k\tau_k(1-\tau_k)^{M_k-1},\\
P_C(k)&=1-P_0(k)-P_E(k)-P_A(k).
\end{align}
$P_E$ sums the disjoint outcomes where exactly one eligible STA transmits;
$P_A$ is the outcome where only the AP transmits. $P_C$ is a collision-event
probability, distinct from conditional collision probability $p_k$.
The boundaries satisfy $P_A(0)=0$ and $P_E(N)=0$.

A successful report increments $k$; AP success decrements it. Thus
\begin{equation}
Q_{k\ell}=\begin{cases}
P_E(k),&\ell=k+1,\\
P_A(k),&\ell=k-1,\\
P_0(k)+P_C(k),&\ell=k,\\
0,&\text{otherwise},
\end{cases}
\end{equation}
with out-of-range transitions omitted. Detailed balance gives
\begin{equation}
\pi_kP_E(k)=\pi_{k+1}P_A(k+1),\qquad 0\leq k<N.
\end{equation}
Writing $w_0=1$, the stationary distribution is
\begin{equation}
w_k=\prod_{j=0}^{k-1}\frac{P_E(j)}{P_A(j+1)},\qquad
\pi_k=\frac{w_k}{\sum_{\ell=0}^{N}w_\ell}.
\end{equation}
Numerically, use $\ell_0=0$ and
$\ell_{k+1}=\ell_k+\log P_E(k)-\log P_A(k+1)$, then normalize
$\exp(\ell_k-\max_j\ell_j)$ to avoid overflow.

Define report and service probabilities per virtual event:
\begin{equation}
q_E=\avg{P_E(k)},\qquad q_A=\avg{P_A(k)}.
\end{equation}
Zero stationary drift implies
\begin{equation}
0=\mathbb E[K_{j+1}-K_j]=q_E-q_A.
\end{equation}
This checks report conservation; it does not validate the approximate state
occupancy against simulation.

\section{Holding times and throughput}
\subsection{Channel duration of each event}
The report event occupies
\begin{equation}
T_E=e/R_C+\delta+d.
\end{equation}
The successful AP exchange has the serialized phases below.
\begin{center}
\begin{tabular}{@{}ll@{}}\toprule
Phase & Duration ($\mu$s)\\\midrule
ICF, LC & $f/R_C$\\
First gap, LC & $s+\delta$\\
Switch up & $u$\\
ICR, HC & $r/R_C$\\
Second gap, HC & $s+\delta$\\
DATA, HC & $(h+L)/R_H$\\
Third gap, HC & $s+\delta$\\
ACK, HC & $a/R_C$\\
Switch down & $v$\\
Final gap, LC & $d+\delta$\\\bottomrule
\end{tabular}
\end{center}
Summing gives
\begin{equation}
T_U=(f+r+a)/R_C+(h+L)/R_H+3s+d+4\delta+u+v.
\end{equation}
HC operation does not accelerate control frames. Since $e=f$, all
contention collisions have fixed frame airtime $t_f=e/R_C$ and
\begin{equation}
T_C=t_f+d+\delta.
\end{equation}
Unequal EBSR/ICF airtimes need collision-type-specific rewards and are
rejected by the current analytical function.

\subsection{Reward normalization}
The conditional mean event duration and its stationary mean are
\begin{equation}
t_k=P_0(k)\sigma+P_E(k)T_E+P_A(k)T_U+P_C(k)T_C,
\qquad \bar t=\avg{t_k}.
\end{equation}
Each AP success earns $L$ payload bits. The event-reward ratio therefore
produces
\begin{align}
S&=\frac{Lq_A}{\bar t}\quad\text{Mbps},\\
\lambda_{del}&=\frac{q_A}{\gamma\bar t}\quad\text{packets/s},\\
C_{packet}&=\frac{\bar t}{1000q_A}\quad\text{ms/delivery}.
\end{align}
$C_{packet}$ is aggregate channel time per delivery, not a tagged packet's
arrival-to-ACK delay.

\section{Collision ratio per transmission attempt}
Let $X_k$ count transmitting contenders in a virtual event. The local
closure makes $X_k$ binomial with parameters $(M_k,\tau_k)$. Thus
\begin{equation}
A_k=\mathbb E[X_k\mid k]=M_k\tau_k.
\end{equation}
A successful event contains exactly one attempt. Removing successes gives
\begin{align}
C_k&=\mathbb E[X_k\ind_{\{X_k\geq2\}}\mid k]\\
&=M_k\tau_k-P_E(k)-P_A(k)\\
&=M_k\tau_k[1-(1-\tau_k)^{M_k-1}]=A_kp_k.
\end{align}
The requested ratio is
\begin{equation}\label{eq:collision}
\boxed{\rho=\frac{\avg{C_k}}{\avg{A_k}}
=\frac{\avg{A_kp_k}}{\avg{A_k}}}.
\end{equation}
It is dimensionless and attempt-weighted. Three colliding senders add three
to each count. Successful EBSR and AP ICF each add one denominator attempt;
protected ICR/DATA/ACK add no separate contention attempts.

For distinction, collision events per second would be
$\sum_k\pi_kP_C(k)/(\gamma\bar t)$, and collision events per delivered
packet would be $\sum_k\pi_kP_C(k)/q_A$. Neither is the current collision
validation metric in Equation~\eqref{eq:collision}.

\section{Time occupancy and channel fractions}
Longer events occupy more elapsed time. Consequently
\begin{equation}
\omega_k=\frac{\pi_kt_k}{\bar t},\qquad
\sum_k\omega_k=1,\qquad \overline K=\sum_k k\omega_k.
\end{equation}
The idle, successful-exchange, and collision time fractions are
\begin{align}
F_0&=\frac{\sigma\sum_k\pi_kP_0(k)}{\bar t},\\
F_U&=\frac{T_Eq_E+T_Uq_A}{\bar t},\\
F_C&=\frac{T_C\sum_k\pi_kP_C(k)}{\bar t}.
\end{align}
By construction $F_0+F_U+F_C=1$. The success fraction includes its gaps
and switches; it is not payload-only airtime.

\section{Detailed all-STA energy rewards}
\subsection{Successful EBSR}
One STA transmits and the other $N-1$ receive on LC. All idle during the
trailing gap, giving
\begin{equation}
E_E=\gamma\left[\frac e{R_C}(P_{L,t}+(N-1)P_{L,r})
+(d+\delta)NP_{L,i}\right].
\end{equation}

\subsection{Successful AP exchange}
Define aggregate powers including all other STAs:
\begin{align}
P_{Ht}&=P_{H,t}+(N-1)P_{L,r},&P_{Hr}&=P_{H,r}+(N-1)P_{L,r},\\
P_{Hi}&=P_{H,i}+(N-1)P_{L,i},&P_{Ss}&=P_{sw}+(N-1)P_{L,i}.
\end{align}
The complete reward, with each term matched to its phase, is
\begin{align}
E_U=\gamma\Bigl[&
\underbrace{\frac f{R_C}NP_{L,r}}_{\text{ICF}}+
\underbrace{(s+\delta)NP_{L,i}}_{\text{first gap}}\notag\\
&+\underbrace{uP_{Ss}}_{\text{switch up}}+
\underbrace{\frac r{R_C}P_{Ht}}_{\text{ICR}}+
\underbrace{(s+\delta)P_{Hi}}_{\text{second gap}}\notag\\
&+\underbrace{\frac{h+L}{R_H}P_{Ht}}_{\text{DATA}}+
\underbrace{(s+\delta)P_{Hi}}_{\text{third gap}}\notag\\
&+\underbrace{\frac a{R_C}P_{Hr}}_{\text{ACK}}+
\underbrace{vP_{Ss}}_{\text{switch down}}+
\underbrace{(d+\delta)NP_{L,i}}_{\text{final gap}}\Bigr].
\end{align}
The AP sends ICF and ACK, but its transmit energy is excluded. Other STAs
use LC receive accounting during HC frames. That is an energy assumption,
not a claim that those STAs decode HC frames.

\subsection{Collision reward without double weighting}
Let $Z_k$ be the number of transmitting STAs, excluding the AP. Its mean
is $n_k\tau_k$. A successful EBSR has one STA sender, and a successful AP
ICF has none. The unconditional collision-restricted sender count is
\begin{equation}
x_k=\mathbb E[Z_k\ind_{\{X_k\geq2\}}\mid k]
=n_k\tau_k-P_E(k).
\end{equation}
The collision-restricted number of nontransmitting STA receivers has mean
$NP_C(k)-x_k$. Hence the already probability-weighted reward is
\begin{equation}
\widetilde E_C(k)=\gamma\left[
t_f\{x_kP_{L,t}+(NP_C(k)-x_k)P_{L,r}\}
+P_C(k)(d+\delta)NP_{L,i}\right].
\end{equation}
Do not multiply by $P_C(k)$ again. This counts all senders without assuming
that a collision always involves exactly two nodes.

\subsection{Mean power and energy per delivery}
The idle-event reward is $E_0=\gamma\sigma NP_{L,i}$. Set
\begin{equation}
E_k=P_0(k)E_0+P_E(k)E_E+P_A(k)E_U+\widetilde E_C(k),
\qquad\bar E=\avg{E_k}.
\end{equation}
Then
\begin{equation}
\boxed{\overline P_{STA}=\frac{\bar E}{N\gamma\bar t}\ \text{W}},
\qquad
\boxed{E_{packet}=1000\frac{\bar E}{q_A}\ \text{mJ}}.
\end{equation}
Packet energy includes the waiting and overhearing energy of all STAs,
normalized by successful deliveries, not only the served STA's own energy.

\section{Intervening attempts after EBSR}
For completed pair $j$, let $B_j$ be the cumulative network contention-attempt
counter just after successful EBSR. Let $G_j$ include the corresponding
successful AP ICF attempt. The simulator measures
\begin{equation}
I_j=G_j-B_j-1,\qquad \widehat I=\frac1D\sum_{j=1}^{D}I_j.
\end{equation}
$D$ counts packets whose ACK completes inside the observation window.
Both endpoints are excluded. Other successful contentions count once;
collisions count each sender. Protected response frames do not contend
separately. With no delivery the mean is undefined (NaN). With one STA it
is zero, since no other contender can intervene after that STA reports.

This metric is simulation-only in the current implementation. Its analytical
prediction would require extra tagged-report/FIFO-position information or a
separate tagged-process analysis. Completed-pair selection can favor shorter
waits in finite windows.

\section{Closed-form checks}
\subsection{Single-STA cycle}
With $N=1$, there is one contender in each state, no collisions, and two
fresh uniform backoffs per delivered packet. Each backoff has mean
$(W-1)/2$ slots. Thus
\begin{equation}
T_{cycle}=(W-1)\sigma+T_E+T_U,\qquad S=L/T_{cycle}.
\end{equation}
Equivalently, $P_E(0)=P_A(1)=\tau=2/(W+1)$ and $\pi_0=\pi_1=1/2$.
The reward formula gives
\begin{equation}
\frac{\bar t}{q_A}
=\frac{2(1-\tau)\sigma+\tau(T_E+T_U)}{\tau}
=(W-1)\sigma+T_E+T_U,
\end{equation}
matching the direct cycle derivation.

\subsection{Uniform-power conservation}
If all LC, HC, and switch powers equal $P$, every STA consumes $P$ at
all times. Then
\begin{align}
E_E&=\gamma NPT_E,&E_U&=\gamma NPT_U,\\
\widetilde E_C(k)&=\gamma NP P_C(k)T_C,&E_k&=\gamma NPt_k.
\end{align}
Therefore $\overline P_{STA}=P$, independently of collisions and switching.

\section{Numerical procedure and validation}
\begin{enumerate}
\item For each $k$, compute the active population and solve
Equation~\eqref{eq:fixed}.
\item Evaluate event probabilities, holding time, and energy rewards.
\item Normalize log weights to obtain the embedded stationary distribution.
\item Check probability normalization and report/service flow balance.
\item Accumulate service, time, attempt, collision, and energy rewards.
\item Divide each reward by its correct denominator; embedded probabilities
must not be substituted for time occupancy.
\end{enumerate}
Run \texttt{py -B tests/check\_models.py} for consistency checks and
\texttt{py -B validate\_dps.py} for fresh saturated comparisons. Five
scenarios, three station counts, and three seeds yield 45 default runs.
Nine further runs examine 10/30/100-second windows at $N=50$.
The numerical results are in
\texttt{results/dps\_validation/comparison.csv} and its accompanying report.

For independent simulated values $y_1,\ldots,y_R$, the script reports
\begin{equation}
\bar y=\frac1R\sum_{r=1}^{R}y_r,\qquad
s_y^2=\frac1{R-1}\sum_{r=1}^{R}(y_r-\bar y)^2,\qquad
\mathrm{SE}=\frac{s_y}{\sqrt R}.
\end{equation}
Relative prediction error is
\begin{equation}
100\frac{y_{analytical}-\bar y}{\bar y}\%,\qquad\bar y\ne0.
\end{equation}
Error bars are one SE, not 95\% confidence intervals. Plot axes begin at
zero; this display choice does not change numerical errors.

The prediction is stationary, whereas simulation starts at $k=0$ with no
pending reports. Counter/stage memory, report identities, and initial
transients are not captured by the reduced closure. Close throughput
agreement does not imply accurate collision dynamics or validated internal
occupancy. Poisson delay/PRR, tagged-report counts, analytical state
occupancy, and PHY/standard correctness are outside this validation scope.

\section*{Attribution}
The DCF fixed-point closure follows the Bianchi decoupling approach
(2000, DOI \texttt{10.1109/49.840210}). The DPS state reduction, holding times,
and reward construction here are this project's approximation, not results
claimed from the original DCF analysis.
\end{document}
